\subsection{Channel models}
We benchmark four channel families.
All channels are implemented on real-valued vectors, but not all channels use the same interpretation of the coordinates.
For AWGN and Rayleigh, the benchmark uses real-valued channel inputs and outputs directly.
In the main benchmark, both channels use dimension \(n=7\).
For SSPA and OptFib, the signal is again stored as a real vector, but adjacent entries are interpreted as paired in-phase and quadrature components.
Thus SSPA uses \(n=8\), corresponding to four I/Q pairs, and OptFib uses \(n=2\), corresponding to one I/Q pair.
Noise variables are sampled independently across real components unless the channel transformation itself introduces coupling.

\textbf{AWGN:}
\begin{equation}
y = x + n,
\qquad
n \sim \mathcal{N}(0,\sigma^2 I).
\end{equation}
This is the standard additive white Gaussian noise reference channel and serves as the simplest near-identity baseline.
In our benchmark, AWGN is treated as a real-valued vector channel rather than as a complex I/Q channel.

\textbf{Rayleigh:}
the transmitted symbol is multiplied elementwise by a random fading amplitude and then corrupted by additive Gaussian noise.
In the implemented benchmark,
\begin{equation}
y = h \odot x + n,
\end{equation}
where \(\odot\) denotes elementwise multiplication,
\begin{equation}
h_k = \frac{1}{\sqrt{2}}\sqrt{u_k^2 + v_k^2},
\qquad
u_k,v_k \overset{\text{i.i.d.}}{\sim} \mathcal{N}(0,1),
\end{equation}
and $n \sim \mathcal{N}(0,\sigma^2 I).$

Each real component is therefore scaled by an independent Rayleigh-distributed magnitude before noise is added.
This channel is multiplicative rather than additive and is correspondingly less aligned with a pure residual perturbation view.
As implemented here, Rayleigh is also a real-valued vector channel rather than a paired-I/Q complex channel.

\textbf{SSPA:}
the input is passed through a smooth solid-state power amplifier nonlinearity before additive noise is applied.
In this case, the real-valued representation is interpreted as paired I/Q components.
Let
\begin{equation}
r = \sqrt{x_I^2 + x_Q^2}
\end{equation}
denote the input amplitude.
The nonlinear amplitude law used in the benchmark is
\begin{equation}
a(r)
=
\frac{G}{\left(1+\left(\frac{Gr}{A_0}\right)^{2p}\right)^{1/(2p)}},
\end{equation}
with default parameters \(p=3\), \(A_0=1.5\), and \(G=5.0\).
The noiseless amplifier output is $\tilde{x} = a(r)\,x,$
and the observed channel output is
\begin{equation}
y = \tilde{x} + n,
\qquad
n \sim \mathcal{N}\!\left(0,\frac{\sigma^2}{2}I\right).
\end{equation}
This channel introduces amplitude-dependent nonlinear distortion while preserving the input direction in the I/Q plane.

\textbf{OptFib:}
We include a simplified memoryless optical-fiber proxy channel following \cite{li2018e2efiber} for the choice of parameters.
As in SSPA, the real-valued state is interpreted as an I/Q pair.
Let the state after propagation step \(k\) be \(x^{(k)} = [x_I^{(k)},x_Q^{(k)}]^\top\).
At each step, the instantaneous nonlinear phase is
\begin{equation}
\phi^{(k)}
=
\gamma
\left(\left(x_I^{(k)}\right)^2 + \left(x_Q^{(k)}\right)^2\right)\frac{L}{K_{\text{step}}}.
\end{equation}
The deterministic Kerr-style rotation is
\begin{equation}
\begin{aligned}
\tilde{x}_I^{(k+1)} &= x_I^{(k)}\cos\phi^{(k)} - x_Q^{(k)}\sin\phi^{(k)}, \\
\tilde{x}_Q^{(k+1)} &= x_I^{(k)}\sin\phi^{(k)} + x_Q^{(k)}\cos\phi^{(k)},
\end{aligned}
\end{equation}
followed by additive Gaussian perturbations at every step:
\begin{equation}
\begin{aligned}
x_I^{(k+1)} &= \tilde{x}_I^{(k+1)} + \sigma_n \xi_I^{(k)}, \\
x_Q^{(k+1)} &= \tilde{x}_Q^{(k+1)} + \sigma_n \xi_Q^{(k)},
\end{aligned}
\qquad
\xi_I^{(k)},\xi_Q^{(k)} \overset{\text{i.i.d.}}{\sim} \mathcal{N}(0,1).
\end{equation}
When that optical-channel parameterization is used,
\begin{equation}
\sigma_n
=
\sqrt{\frac{10^{(P_n-30)/10}}{2K_{\text{step}}}},
\end{equation}
with \(P_n\) measured in dBm.
